It is important to note that our methodology involves multiple degrees of freedom, as this study is inherently exploratory in nature. First, our cross-sectional design limits the ability to infer individual developmental trajectories. Longitudinal studies tracking the same individuals over time would provide more definitive evidence of how SFC evolves within subjects. Second, while we utilized advanced imaging techniques to estimate structural and functional connectivity, these methods have inherent limitations in spatial and temporal resolution. Future studies employing higher-resolution imaging modalities could yield more precise characterizations of SFC dynamics. Third, our focus on typically developing children precludes examination of how atypical development or neurodevelopmental disorders may impact SFC maturation. Investigating clinical populations could elucidate how deviations from normative trajectories relate to cognitive and behavioral outcomes. Finally, while we identified associations between SFC, myelination, and excitation-inhibition balance, the causal mechanisms underlying these relationships remain unclear. Future research integrating multimodal imaging with molecular and cellular approaches could help unravel the biological processes driving SFC development.

